% ============================================================
%  §8  Pathologies and Limits
% ============================================================

\section{Pathologies and Limits}
\label{sec:pathologies}

\subsection{Deliberate Exclusions}

The present framework deliberately excludes the following:

\textbf{Semantic theory.}
No claims are made about reference, truth conditions, or
interpretative content.
Meaning, where it appears, is an emergent label applied to stable
trajectories; it is not a primitive of the theory.

\textbf{Phenomenology and consciousness.}
The framework does not generate or explain subjective experience.
Conscious access is treated as a downstream phenomenon; its
mechanistic basis is the subject of a companion paper on
Consciousness Theory.

\textbf{Neurobiological implementation.}
The theory describes structural and dynamical conditions at a level
of abstraction that is independent of specific neural architectures.
Correspondence to neural implementation is deferred to future work.

\subsection{Failure Modes}

\textbf{Overclosure (too-deep wells).}
When the effective $\sigma_{\mathrm{barrier}}$ grows large relative to
$\sigma_{\mathrm{drive}}$, existing wells deepen to the point where
new input can no longer produce topological change.
The landscape hardens; thought becomes repetitive and loses
adaptability.
This corresponds phenomenologically to rigidity, dogmatism, or
entrenched belief.

\textbf{Underclosure (too-shallow wells).}
When $\sigma_{\mathrm{dissip}}$ dominates, existing wells shallow
faster than new ones can form.
The landscape remains flat; thought becomes unstable and fails to
accumulate.
This corresponds to the failure of concepts to stabilize, preventing
the formation of durable understanding.

\textbf{Synchronization breakdown (social scale).}
When linguistic tokens fail to produce correlated perturbations across
agents --- due to divergent prior landscapes, insufficient shared
exposure, or deliberate manipulation --- the social landscape
fragments.
The result is the loss of shared conceptual terrain, which the
framework identifies as the structural basis of communication failure,
cultural conflict, and epistemic polarization.

\subsection{Scope Limits}

The theory establishes structural and dynamical conditions; it does
not establish:
\begin{itemize}
  \item Quantitative predictions about specific cognitive systems
        (this requires concrete specification of $K$, $D$, $\mu$,
        $\alpha$, $\Gamma$, $\kappa_B$).
  \item Correspondence conditions between $\Phi$-geometry and
        observable neural or behavioral data.
  \item A complete theory of language acquisition, semantic change,
        or collective intelligence.
\end{itemize}
These remain as future work, on the scaffold established here.
